\documentclass[12pt]{article}
\usepackage{amsmath}
\usepackage{amsfonts}
\usepackage{amssymb}
\usepackage{geometry}
\usepackage{tikz}
\geometry{a4paper, margin=1in}
\pagenumbering{gobble}

\title{02YMECH - Homework 4}
\date{Assigned for the week of Oct 13, 2025}
\author{}
\begin{document}
\maketitle

\section*{Questions}

\begin{enumerate}
    \item The position of a particle hooked to a spring is given by the equation \( x(t)=A \cos(\omega_0 t + \phi) \), where \( A \), \( \omega_0 \), and \( \phi \) are constants. If the particle's initial position is $x_0$ and initial velocity is $v_0$, how long will it take for the particle to reach its maximum displacement for the first time?
    \item What is the average kinetic energy and average potential energy of a harmonic oscillator over one complete cycle? Express your answers in terms of the total energy \( E \) of the system.(Consider the time integrals of kinetic and potential energy over one period.)
    \item A simple pendulum consists of a bob of mass $m$ attached to a light, inextensible string of length $L$, staying in its equilibrium position. The support point of the pendulum is suddenly accelerated horizontally with constant acceleration $a$ (e.g., the ceiling begins to move to the right). 
    \begin{enumerate}
        \item Determine the new equilibrium angle \( \theta \) that the pendulum makes with the vertical.
        \item Find the period of small oscillations about this new equilibrium position.
    \end{enumerate}
    \item The position of a damped harmonic oscillator is given by the equation \( x(t) = A e^{-\beta t} \cos(\omega_1 t + \phi) \), where \( A \), \( \beta \), \( \omega_1 \), and \( \phi \) are constants. How much time will it take for the amplitude of the oscillations to decrease to half of its initial value?

\end{enumerate}
\end{document}
